% SGA TRANSLATION PROJECT
% Volume: SGA 4
% Expose: IV, Topoi
% Source: Orgogozo/Laszlo TeX snapshot, SGA4 orgogozo_tex_sources/04/04.tex
% Range translated in this working draft section: Expose IV, Sections 10--14; source lines 8246--9852.
% Package policy: fresh/new material only, not cumulative, to keep downloads small.
% Status: draft translation, compile-checked, not proofed line-by-line against scan.

\documentclass[11pt]{article}
\usepackage[T1]{fontenc}
\usepackage[utf8]{inputenc}
\usepackage[english]{babel}
\usepackage{lmodern}
\usepackage{microtype}
\usepackage{amsmath,amssymb,amsthm,mathtools}
\usepackage{mathrsfs}
\usepackage{enumitem}
\usepackage{geometry}
\geometry{margin=1.10in}
\usepackage[all]{xy}
\usepackage{hyperref}
\hypersetup{colorlinks=true,linkcolor=blue,citecolor=blue,urlcolor=blue,hypertexnames=false}

\newcommand{\U}{\mathcal{U}}
\newcommand{\V}{\mathcal{V}}
\newcommand{\C}{\mathcal{C}}
\newcommand{\E}{\mathcal{E}}
\newcommand{\F}{\mathcal{F}}
\newcommand{\Oo}{\mathcal{O}}
\newcommand{\ZZ}{\mathbb{Z}}
\newcommand{\Set}{\mathbf{Set}}
\newcommand{\Ab}{\mathbf{Ab}}
\newcommand{\Top}{\operatorname{Top}}
\newcommand{\Hom}{\operatorname{Hom}}
\newcommand{\ob}{\operatorname{ob}}
\newcommand{\id}{\operatorname{id}}
\newcommand{\Ker}{\operatorname{Ker}}
\newcommand{\Spec}{\operatorname{Spec}}
\newcommand{\Sup}{\operatorname{Sup}}
\newcommand{\Res}{\operatorname{Res}}
\newcommand{\Ppoint}{\operatorname{Pt}}
\newcommand{\Pptgeom}{\operatorname{Pt}_{\mathrm{geom}}}
\newcommand{\Hhomtop}{\mathcal{H}om_{\mathrm{top}}}
\newcommand{\Hhomtopan}{\mathcal{H}om_{\mathrm{top,ann}}}
\newcommand{\Hhomtoplocan}{\mathcal{H}om_{\mathrm{top,loc.ann}}}
\newcommand{\scrHom}{\mathscr{H}\!om}
\newcommand{\ie}{i.e.}
\newcommand{\cf}{cf.}
\newcommand{\resp}{respectively}
\newcommand{\loccit}{loc. cit.}
\newcommand{\quav}{\quad\text{and}\quad}
\newcommand{\wh}[1]{\widehat{#1}}
\newcommand{\tw}[1]{\widetilde{#1}}
\newcommand{\ad}{\operatorname{ad}}
\newcommand{\ab}{\mathrm{ab}}

\newenvironment{statement}[1]{\par\medskip\noindent\textbf{#1.}\ }{\par\medskip}
\newenvironment{remarktext}[1]{\par\medskip\noindent\textbf{#1.}\ }{\par\medskip}
\newenvironment{prooftext}{\par\noindent\emph{Proof.}\ }{\par\medskip}

\begin{document}

\begin{center}
{\Large \textbf{SGA 4, Expos\'e IV. Topoi}}\\[0.35em]
{\large Draft translation, Sections 10--14}\\[0.35em]
A. Grothendieck and J.-L. Verdier
\end{center}

\bigskip


\setcounter{section}{9}
\section{Sheaves of morphisms}

\begin{statement}{Proposition 10.1}
Let $E$ be a topos, and let $X$ and $Y$ be two objects of $E$. The functor
\[
Z\longmapsto \Hom_E(Z\times X,Y)\simeq \Hom_{E_{/Z}}(X_Z,Y_Z)
\]
is representable.
\end{statement}

Indeed, inductive limits are universal in $E$ (II, 4.3). Hence the functor $Z\mapsto Z\times X$ commutes with inductive limits. It follows that $Z\mapsto \Hom_E(Z\times X,Y)$ carries inductive limits in the argument $Z$ into projective limits. It is therefore representable (IV, 1.4 and 1.2, equivalently 1.2.1).

\subsection*{10.2}
The object representing the functor $Z\mapsto \Hom_E(Z\times X,Y)$ is denoted
\[
\scrHom_E(X,Y),
\]
or simply $\scrHom(X,Y)$, and is called the \emph{sheaf of morphisms from $X$ to $Y$}. It is bifunctorial in $X$ and $Y$. Thus one has a trifunctorial isomorphism
\begin{equation}
\Hom_E(Z,\scrHom_E(X,Y))\simeq \Hom_E(Z\times X,Y).\tag{10.2.1}
\end{equation}
It follows from (10.2.1) that the bifunctor $(X,Y)\mapsto \scrHom(X,Y)$ sends inductive limits in the argument $X$ (respectively projective limits in the argument $Y$) to projective limits in $E$.

\begin{statement}{Proposition 10.3}
Let $v:E\to E'$ be a geometric morphism. For a variable object $X$ of $E$ and a variable object $Y$ of $E'$, there is a bifunctorial isomorphism
\begin{equation}
v_*\scrHom_E(v^*Y,X)\simeq \scrHom_{E'}(Y,v_*X).\tag{10.3.1}
\end{equation}
\end{statement}

\begin{prooftext}
For every object $Z$ of $E'$ there is a chain of isomorphisms, functorial in all arguments:
\begin{align*}
\Hom_{E'}(Z,v_*\scrHom_E(v^*Y,X))
&\simeq \Hom_E(v^*Z,\scrHom_E(v^*Y,X)) &&\text{(adjunction)}\\
&\simeq \Hom_E(v^*Z\times v^*Y,X) &&\text{by (10.2.1)}\\
&\simeq \Hom_E(v^*(Z\times Y),X) &&\text{since $v^*$ is left exact}\\
&\simeq \Hom_{E'}(Z\times Y,v_*X) &&\text{(adjunction)}\\
&\simeq \Hom_{E'}(Z,\scrHom_{E'}(Y,v_*X)) &&\text{by (10.2.1).}
\end{align*}
Thus the two sides of (10.3.1) represent isomorphic functors.
\end{prooftext}

\begin{statement}{Corollary 10.4}
Let $C$ be a $\U$-site, $X$ a $\U$-presheaf on $C$, $Y$ a $\U$-sheaf on $C$, and $\V$ a universe containing $\U$ such that $C$ is $\V$-small. Let $\wh C_{\V}$ be the topos of $\V$-presheaves on $C$, and let $C^{\sim}_{\U}$ be the topos of $\U$-sheaves on $C$. Then the presheaf
\[
\scrHom_{\wh C_{\V}}(X,Y)
\]
is a $\U$-sheaf. If $X\to \underline{a}X$ is the canonical morphism from $X$ to its associated sheaf, the corresponding morphism
\[
\scrHom_{\wh C_{\V}}(\underline{a}X,Y)\longrightarrow \scrHom_{\wh C_{\V}}(X,Y)
\]
is an isomorphism. Moreover there is a canonical isomorphism
\[
\scrHom_{\wh C_{\V}}(\underline{a}X,Y)\simeq \scrHom_{C^{\sim}_{\U}}(\underline{a}X,Y).
\]
\end{statement}

\subsubsection*{10.4.1}
First suppose that $C$ is a small site and that $\V=\U$. There is then a geometric morphism
\[
C^{\sim}_{\U}\longrightarrow \wh C_{\U}
\]
(IV, 4.8), and (10.3.1) gives the assertions in this case. To pass from this to the general case, first observe that the associated-sheaf functor is independent of the universe (II, 3.6), and that the topos of $\V$-sheaves on $C$ is equivalent to the topos of $\V$-sheaves on $C^{\sim}_{\U}$ (III, 4 and IV, 1). It remains only to prove the following lemma.

\begin{statement}{Lemma 10.4.2}
Let $E$ be a $\U$-topos, let $\V$ be a universe containing $\U$, let $E^{\sim}_{\V}$ be the topos of $\V$-sheaves on $E$, and let $X,Y$ be two objects of $E$. There is a canonical isomorphism
\[
\scrHom_E(X,Y)\simeq \scrHom_{E^{\sim}_{\V}}(X,Y).
\]
\end{statement}

\subsubsection*{10.4.3}
The functor $\epsilon_E:E\to E^{\sim}_{\V}$ is fully faithful and left exact. Therefore, for every object $Z$ of $E$, there is an isomorphism
\[
\Hom_{E^{\sim}_{\V}}(\epsilon_EZ,\epsilon_E\scrHom_E(X,Y))\simeq
\Hom_{E^{\sim}_{\V}}(\epsilon_EZ\times\epsilon_EX,\epsilon_EY).
\]
But every object of $E^{\sim}_{\V}$ is an inductive limit of objects coming from $E$; by (10.2.1) this proves the assertion.

\subsection*{10.5}
Let $v:E\to E'$ be a geometric morphism. We shall define four bifunctorial morphisms:
\begin{align*}
\Phi_v&:v^*\scrHom(X,Y)\longrightarrow \scrHom(v^*X,v^*Y),
& &X,Y\in\ob E',\\
\Psi_v&:v_*\scrHom(X,Y)\longrightarrow \scrHom(v_*X,v_*Y),
& &X,Y\in\ob E,\\
\Xi_v&:\scrHom(v_!X,Y)\longrightarrow v_*\scrHom(X,v^*Y),
& &X\in\ob E,\ Y\in\ob E',\\
\Lambda_v&:v_!(X\times v^*Y)\longrightarrow v_!(X)\times Y,
& &X\in\ob E,\ Y\in\ob E'.
\end{align*}
In the definitions of $\Xi_v$ and $\Lambda_v$ one assumes that the inverse-image functor $v^*$ has a left adjoint $v_!$.

\subsubsection*{10.5.1. Definition of $\Psi_v$}
Let $X$ be an object of $E$. The adjunction morphism $v^*v_*X\to X$ gives, for every object $Z$ of $E'$, a bifunctorial morphism
\[
v^*(Z\times v_*X)\longrightarrow v^*Z\times X.
\]
Thus, for every object $Y$ of $E$, one has a trifunctorial morphism
\[
\Hom_E(v^*Z\times X,Y)\longrightarrow \Hom_E(v^*(Z\times v_*X),Y).
\]
By adjunction this gives a trifunctorial morphism
\[
\Hom_{E'}(Z,v_*\scrHom(X,Y))\longrightarrow
\Hom_{E'}(Z,\scrHom(v_*X,v_*Y)),
\]
and hence the morphism $\Psi_v$.

\subsubsection*{10.5.2. Definition of $\Lambda_v$}
For every object $X$ of $E$ there is an adjunction morphism $X\to v^*v_!X$. Therefore, for every object $Y$ of $E'$, one obtains a bifunctorial morphism
\[
X\times v^*Y\longrightarrow v^*v_!(X)\times v^*Y\simeq v^*(v_!(X)\times Y),
\]
and hence, by adjunction, the morphism $\Lambda_v$.

\subsubsection*{10.5.3. Definition of $\Xi_v$}
Let $X$ be an object of $E$, and let $Z,Y$ be objects of $E'$. The morphism
\[
\Lambda_v:v_!(X\times v^*Z)\to v_!(X)\times Z
\]
gives a trifunctorial morphism
\[
\Hom_{E'}(v_!(X)\times Z,Y)\longrightarrow
\Hom_{E'}(v_!(X\times v^*Z),Y).
\]
By adjunction this gives a trifunctorial morphism
\[
\Hom_{E'}(Z,\scrHom(v_!X,Y))\longrightarrow
\Hom_{E'}(Z,v_*\scrHom(X,v^*Y)),
\]
and hence the morphism $\Xi_v$.

\subsubsection*{10.5.4. Definition of $\Phi_v$}
The trifunctional morphism of 10.5.3,
\[
\Hom_{E'}(v_!Z\times X,Y)\longrightarrow \Hom_{E'}(v_!(Z\times v^*X),Y),
\]
where $Z$ is an object of $E$ and $X,Y$ are objects of $E'$, gives by adjunction a trifunctional morphism
\[
\Hom_E(Z,v^*\scrHom(X,Y))\longrightarrow
\Hom_E(Z,\scrHom(v^*X,v^*Y)),
\]
and hence one definition of $\Phi_v$.

There is a second way to define $\Phi_v$. Let $X,Y,Z$ be three objects of $E'$. The functor $v^*:E'\to E$ gives a trifunctional morphism
\[
\Hom_{E'}(Z\times X,Y)\longrightarrow
\Hom_E(v^*Z\times v^*X,v^*Y).
\]
By adjunction this yields
\[
\Hom_{E'}(Z,\scrHom(X,Y))\longrightarrow
\Hom_{E'}(Z,v_*\scrHom(v^*X,v^*Y)).
\]
Thus one obtains a bifunctorial morphism
\[
\scrHom(X,Y)\longrightarrow v_*\scrHom(v^*X,v^*Y),
\]
and, again by adjunction, a morphism which the reader may verify is the same as the morphism $\Phi_v$ defined above.

\begin{statement}{Proposition 10.6}
Let $v:E\to E'$ be a geometric morphism.
\begin{enumerate}[label=\arabic*)]
\item The morphism $\Psi_v$ of 10.5.1 is an isomorphism for every object $Y$ of $E$ if and only if the adjunction morphism $v^*v_*X\to X$ is an isomorphism. In particular, $\Psi_v$ is an isomorphism for every pair $(X,Y)$ if and only if $v_*:E\to E'$ is fully faithful, \ie\ if and only if $v$ is an embedding of topoi.
\item The following conditions are equivalent:
\begin{enumerate}[label=(\roman*)]
\item for every pair $(X,Y)$ of objects of $E'$, the morphism $\Phi_v$ of 10.5.4 is an isomorphism;
\item the functor $v^*$ has a left adjoint $v_!$ and, for every object $X$ of $E$ and every object $Y$ of $E'$, the morphism $\Xi_v$ of 10.5.3 is an isomorphism;
\item the functor $v^*$ has a left adjoint $v_!$ and, for every object $X$ of $E$ and every object $Y$ of $E'$, the morphism $\Lambda_v$ of 10.5.2 is an isomorphism.
\end{enumerate}
\end{enumerate}
\end{statement}

\subsubsection*{10.6.1}
The first assertion follows immediately from 10.5.1. It also follows immediately from 10.5.2, 10.5.3, and 10.5.4 that (iii) is equivalent to (ii), and that these are equivalent to (i) together with the existence of a left adjoint $v_!$ to $v^*$. It remains to prove that (i) implies that $v^*$ has a left adjoint, \ie\ by IV, 1.8, that (i) implies that $v^*$ commutes with small products. Let $e'$ be the final object of $E'$ and let $I$ be a small set. Since $v^*$ is left exact, $v^*(e')$ is final in $E$; since $v^*$ commutes with inductive limits, $v^*(\coprod_I e')\simeq\coprod_Iv^*(e')$. For every object $Y$ of $E'$ one has
\[
\scrHom(\coprod_Ie',Y)\simeq \prod_I\scrHom(e',Y)\simeq \prod_IY
\]
and
\[
\scrHom(\coprod_Iv^*(e'),v^*Y)\simeq \prod_I\scrHom(v^*(e'),v^*Y)\simeq \prod_Iv^*Y.
\]
The functorial morphism $\Phi_v$ therefore induces an isomorphism
\[
v^*\prod_IY\simeq \prod_Iv^*Y,
\]
and one checks that this is the canonical isomorphism.

\begin{statement}{Corollary 10.7}
Let $E$ be a topos, let $Z$ be an object of $E$, let $j_Z:E_{/Z}\to E$ be the localization morphism, and let $X,Y$ be two objects of $E$.
\begin{enumerate}[label=\arabic*)]
\item There is a bifunctorial isomorphism
\[
j_Z^*\scrHom(X,Y)\simeq \scrHom(j_Z^*X,j_Z^*Y).
\]
\item If $X'$ is an object of $E_{/Z}$, there is a bifunctorial isomorphism
\[
\scrHom_E(j_{Z!}X',Y)\simeq j_{Z*}\scrHom_{E_{/Z}}(X',j_Z^*Y).
\]
\item If $Y'$ is an object of $E_{/Z}$ and $Z$ is an open of $E$, then there is a bifunctorial isomorphism
\[
j_{Z*}\scrHom(X',Y')\simeq \scrHom(j_{Z*}X',j_{Z*}Y').
\]
\end{enumerate}
\end{statement}

Assertions 1) and 2) are proved by observing that $\Lambda_{j_Z}$ is an isomorphism. For 3) it is enough to observe that $j_{Z*}$ is fully faithful, because $j_{Z!}$ is fully faithful (I, 5.7 a)).

\begin{statement}{Corollary 10.8}
Let $E$ be a topos, let $X,Y,Z$ be three objects of $E$, and let $j_Z:E_{/Z}\to E$ be the localization morphism.
\begin{enumerate}[label=\arabic*)]
\item There is a canonical isomorphism
\[
j_{Z*}j_Z^*Y\simeq \scrHom(Z,Y).
\]
\item There is a canonical isomorphism
\[
\Hom_E(Z,\scrHom(X,Y))\simeq \Hom_{E_{/Z}}(j_Z^*X,j_Z^*Y).
\]
\end{enumerate}
\end{statement}

Let $e_Z$ be the final object of $E_{/Z}$, \ie\ the object $\id_Z:Z\to Z$. From (10.2.1) one has a canonical isomorphism
\[
\scrHom_{E_{/Z}}(e_Z,j_Z^*Y)\simeq j_Z^*Y.
\]
By Corollary 10.7, 2), this gives an isomorphism
\[
j_{Z*}j_Z^*Y\simeq \scrHom(j_{Z!}e_Z,Y).
\]
Since $j_{Z!}e_Z=Z$, this proves 1). To prove 2), Corollary 10.7 gives a canonical isomorphism
\[
\Hom_{E_{/Z}}(e_Z,j_Z^*\scrHom(X,Y))\simeq
\Hom_{E_{/Z}}(j_Z^*X,j_Z^*Y).
\]
Adjunction on the first member then gives the asserted isomorphism.

\section{Ringed topoi, localization in ringed topoi}

\subsubsection*{11.1.1}
Let $\U$ be a universe. A $\U$-\emph{ringed topos} is a pair $(E,A)$, where $E$ is a $\U$-topos and $A$ is an object endowed with a ring structure. A $\U$-\emph{ringed site} is a pair $(C,A)$, where $C$ is a $\U$-site and $A$ is a $\U$-sheaf of rings on $C$. The universe will not be mentioned when there is no risk of ambiguity. A ringed site $(C,A)$ has associated ringed topos $(C^{\sim},A)$. A ringed topos $(E,A)$ has associated ringed site consisting of the site $E$ and the sheaf of rings represented by $A$.

\subsubsection*{11.1.2}
Let $(E,A)$ be a ringed topos. Write ${}_AE$ (respectively $E_A$) for the category of sheaves of unital left (respectively right) $A$-modules; we shall also write ``$A$-modules''. The category ${}_AE$ (respectively $E_A$) is abelian (II, 6.7). If $M$ and $N$ are two sheaves of left (respectively right) $A$-modules, the abelian group of $A$-linear morphisms from $M$ to $N$ is denoted $\Hom_A(M,N)$.

\subsubsection*{11.1.3}
Let $A,B,C$ be three rings in a topos $E$, let $M$ be a sheaf of $A$-$B$-bimodules, and let $N$ be a sheaf of $A$-$C$-bimodules on the left. Let $e$ be a final object of $E$ and put $\Hom(e,B)=\Gamma(B)$ and $\Hom(e,C)=\Gamma(C)$. The $A$-$B$-bimodule structure on $M$ gives a ring homomorphism from $\Gamma(B)^{\circ}$ to the endomorphism ring of the $A$-module $M$, and similarly the $A$-$C$-bimodule structure on $N$ gives a ring homomorphism from $\Gamma(C)^{\circ}$ to the endomorphism ring of the $A$-module $N$. By functoriality one obtains a $\Gamma(B)$-$\Gamma(C)$-bimodule structure on the abelian group $\Hom_A(M,N)$. In particular, when $A$ is a sheaf of commutative rings, $\Hom_A(M,N)$ is canonically a $\Gamma(A)$-module.

\subsubsection*{11.1.4}
Let $E$ be a topos, let $A$ and $B$ be two rings in $E$, let $u:A\to B$ be a morphism of sheaves of rings, and let $M$ be a $B$-module, say on the left. The sheaf $M$ may be regarded as a sheaf of $A$-modules via $u$: for every object $X$ of $E$, the $A(X)$-module structure on $M(X)$ is deduced from its $B(X)$-module structure and from $u(X):A(X)\to B(X)$ by restriction of scalars. Thus one obtains a functor
\[
\Res(u):{}_BE\longrightarrow {}_AE,
\]
called \emph{restriction of scalars by $u$}.

\subsubsection*{11.1.5}
In particular, when $A=\ZZ$ is the constant sheaf $\underline{\ZZ}$ (the sheaf associated with the constant presheaf $\ZZ$), and when $u:\underline{\ZZ}\to B$ is the unique canonical morphism, one obtains a functor from ${}_BE$ to the category ${}_{\underline{\ZZ}}E$, which is none other than the category of sheaves of abelian groups, denoted $E_{\Ab}$. This functor is called the \emph{underlying abelian sheaf functor}.

\begin{statement}{Proposition 11.1.6}
Restriction of scalars by $u$ commutes with inductive and projective limits. It is conservative.
\end{statement}

The proposition is true for the ordinary restriction-of-scalars functor for modules, \ie\ when $E$ is the punctual topos. It is therefore true when $E$ is the topos of presheaves on a small site. The description of inductive and projective limits by means of the associated-sheaf functor (II, 6.4) then implies the proposition in the general case.

\subsubsection*{11.2.1}
Let $(E,A)$ be a ringed topos, let $X$ be an object of $E$, and let $j_X:E_{/X}\to E$ be the localization morphism. By III, 1.7, or IV, 3.1.2, the sheaf $j_X^*A$ has a canonical ring structure. It is usually denoted $A|X$, or abusively again by $A$. Unless otherwise stated, the topos $E_{/X}$ will be ringed by $A|X$. The sheaf $j_{X*}j_X^*A$ has a canonical ring structure, and the adjunction morphism
\[
A\longrightarrow j_{X*}j_X^*A
\]
is a morphism of sheaves of rings.

\subsubsection*{11.2.2}
Let in addition $M$ be an $A$-module, say on the left. The sheaf $j_X^*M$ has an $A|X$-module structure (III, 1.7, or IV, 3.1.2). This gives a functor
\[
j_X^*:{}_AE\longrightarrow {}_{A|X}E_{/X},
\]
called restriction to $E_{/X}$, or simply restriction to $X$. The functor $j_X^*$ commutes with inductive and projective limits, hence in particular is exact. If $N$ is an $A|X$-module, then $j_{X*}N$ is a sheaf of $j_{X*}(A|X)$-modules (III, 1.7 or IV, 3.1.2), and therefore, by restriction of scalars through the adjunction morphism $A\to j_{X*}(A|X)$ (11.1.4), it is an $A$-module, again denoted $j_{X*}N$. We have thus defined a functor
\[
j_{X*}:{}_{A|X}E_{/X}\longrightarrow {}_AE,
\]
which commutes with projective limits (11.1.6 and III, 1.7). For every $A$-module $M$, the adjunction morphism $M\to j_{X*}j_X^*M$ is a morphism of $A$-modules. For every $A|X$-module $N$ and every $A$-module $M$, this adjunction morphism defines a morphism, bifunctorial in $M$ and $N$,
\begin{equation}
\Hom_{A|X}(j_X^*M,N)\longrightarrow \Hom_A(M,j_{X*}N).\tag{11.2.2.1}
\end{equation}
This morphism is an isomorphism; equivalently, the functors $j_X^*$ and $j_{X*}$ for $A$-modules are adjoint (III, 1.7).

\subsubsection*{11.2.3}
The functors $j_X^*$ and $j_{X*}$ for modules commute with the underlying-set functor (III, 1.7, 4). They therefore commute with the restriction-of-scalars functors, and in particular with the underlying abelian sheaf functor.

\begin{statement}{Proposition 11.3.1}
Let $(E,A)$ be a ringed topos and let $X$ be an object of $E$. The functor
\[
j_X^*:{}_AE\longrightarrow {}_{A|X}E_{/X}
\]
has a left adjoint
\[
j_{X!}:{}_{A|X}E_{/X}\longrightarrow {}_AE,
\]
called \emph{extension by zero}. The extension-by-zero functor is exact and faithful and commutes with inductive limits. The functors $j_{X!}$ for modules commute with restriction of scalars and, in particular, with the underlying abelian sheaf functor.
\end{statement}

The existence of $j_{X!}$ follows from III, 1.7. Since $j_{X!}$ is a left adjoint, it commutes with inductive limits (I, 2). To prove the other assertions, first suppose that $E$ is the topos of presheaves of sets on a small category $C$ containing the object $X$. Then $E_{/X}$ is equivalent to the category of presheaves on $C_{/X}$, and, modulo this equivalence, the functor $j_X^*:E\to E_{/X}$ is just composition with the forgetful functor $C_{/X}\to C$ (I, 5.11). It follows immediately from the explicit construction of $j_{X!}$ (I, 5.1) that, for every $A|X$-module $N$ and every object $Y$ of $C$, one has
\[
j_{X!}N(Y)=\bigoplus_{u\in\Hom_C(Y,X)}N(u).
\]
This gives the exactness of $j_{X!}$ and the fact that extension by zero commutes with restriction of scalars in this case. In the general case one may suppose that $E$ is the topos of sheaves on a small site $C$ and that $X$ comes from an object of $C$ (IV, 1.2.1). Then $E_{/X}$ is equivalent to the topos of sheaves on $C_{/X}$ with its induced topology (III, 5.4), and the geometric morphism $j_X:E_{/X}\to E$ comes from the forgetful functor $C_{/X}\to C$, which is continuous and cocontinuous (III, 5.2). It follows from III, 1.7 that extension by zero for sheaves is obtained by composing extension by zero for presheaves with the associated-sheaf functor. This proves exactness and commutation with restriction of scalars.

To prove faithfulness, it is equivalent to show that, for every $A|X$-module $N$, the adjunction morphism
\[
\ad_N:N\longrightarrow j_X^*j_{X!}N
\]
is a monomorphism. Let $j_{\widehat X!}$ denote extension by zero for presheaves and let $\underline a$ denote the associated-sheaf functor. There is a commutative diagram
\[
\xymatrix@C=2.6cm@R=1.3cm{
N \ar[r]^{\ad_N} \ar[rd]_{\ad_{\widehat N}} & j_X^*\underline a\,j_{\widehat X!}N\\
& \underline a\,j_X^*j_{\widehat X!}N \ar[u]
}
\]
where $\ad_{\widehat N}$ is the adjunction morphism for presheaves. The morphism $\ad_{\widehat N}$ is a monomorphism by what has just been proved, hence so is $\underline a(\ad_{\widehat N})$. Moreover, since the forgetful functor $C_{/X}\to C$ is continuous and cocontinuous, the canonical morphism $\underline a\,j_X^*\to j_X^*\underline a$ is an isomorphism (III, 2.3). Hence $\ad_N$ is a monomorphism.

\begin{remarktext}{Remark 11.3.2}
Let $N$ be an $A|X$-module. The underlying sheaf of sets of $j_{X!}N$ is not, in general, isomorphic to the sheaf obtained by extending the underlying sheaf of sets of $N$ by the empty set (III, 5.3 and IV, 5.2). One should therefore take care not to confuse the functor \emph{extension by zero}, denoted
\[
j_{X!}:{}_{A|X}E_{/X}\to {}_AE
\]
in 11.3.1, with the functor extension by the empty set, also denoted $j_{X!}:E_{/X}\to E$ in III, 5.3 and IV, 5.2. In most cases encountered in practice, the abuse of notation just mentioned causes no confusion. When confusion is nevertheless possible, we shall write $j^{\ab}_{X!}$ and $j^{\Set}_{X!}$ for extension by zero and extension by the empty set, respectively.
\end{remarktext}

\begin{statement}{Proposition 11.3.3}
Let $(E,A)$ be a ringed topos and let $X$ be an object of $E$. The $A$-module $j_{X!}(A|X)$, usually denoted $A_X$ or $A_{X,E}$, is the free $A$-module generated by $X$ (II, 6.5): for every $A$-module $M$ there is a canonical isomorphism, functorial in $M$,
\[
\Hom_E(X,M)\xleftarrow{\sim}\Hom_A(A_X,M).
\]
If $(X_i)_{i\in I}$ is a topologically generating family of $E$ (II, 3.0.1), then $(A_{X_i})_{i\in I}$ is a generating family of the category of $A$-modules.
\end{statement}

Let $e_X$ be the final object of $E_{/X}$, namely $X\xrightarrow{\id}X$. There is an isomorphism
\[
\Hom_{A|X}(A|X,j_X^*M)\xrightarrow{\sim}\Hom_{E_{/X}}(e_X,j_X^*M)
\]
induced by the unit section $e_X\to A|X$. By adjunction this gives an isomorphism
\[
\Hom_A(j^{\ab}_{X!}(A|X),M)\xrightarrow{\sim}\Hom_E(j^{\Set}_{X!}(e_X),M).
\]
Since $j^{\Set}_{X!}(e_X)=X$, one obtains the announced isomorphism. The final assertion follows from II, 6.6.

\begin{remarktext}{Remark 11.3.4}
Let $A$ and $B$ be two sheaves of rings on a topos $E$, let $X$ be an object of $E$, and let $N$ be an $A|X$-$B|X$-bimodule. The abelian sheaf obtained by extending $N$ by zero has a canonical $A$-$B$-bimodule structure, as follows immediately from its explicit description (11.3.1). In particular the free $A$-module $A_X$ generated by $X$ is an $A$-bimodule.
\end{remarktext}

\begin{remarktext}{Exercise 11.3.5}
Let $E$ be a topos, let $X$ be an object of $E$, let $x:P\to E$ be a point of $E$, and let $(x_i)_{i\in I}$ be the family of points of $E_{/X}$ over $x$ (in bijective correspondence with the fiber $X_x$, by (6.7.2)). Show that for every abelian sheaf $M$ on $E_{/X}$, the fiber $(j_{X!}M)_x$ is canonically isomorphic to $\bigoplus_i M_{x_i}$.
\end{remarktext}

\section{Operations on modules}

\begin{statement}{Proposition 12.1}
Let $(E,A)$ be a ringed topos, and let $M,N$ be two left (respectively right) $A$-modules. The functor on $E$ which sends an object $X$ of $E$ to the abelian group
\[
\Hom_A(M,\scrHom_E(X,N))
\]
is representable by an abelian sheaf denoted $\scrHom_A(M,N)$, or sometimes $\scrHom(M,N)$ when no confusion can result. For every object $X$ of $E$ there is a canonical isomorphism
\begin{equation}
\scrHom_A(M,N)(X)\simeq \Hom_{A|X}(j_X^*M,j_X^*N).\tag{12.1.1}
\end{equation}
\end{statement}

There is a canonical isomorphism $\scrHom_E(X,N)=j_{X*}j_X^*N$ (10.8); hence $\scrHom_E(X,N)$ is functorially in $X$ a left (respectively right) $A$-module. The functor $X\mapsto\scrHom_E(X,N)$ transforms inductive limits in $X$ into projective limits of $A$-modules (10.2 and 11.2.3), and therefore the functor
\[
X\longmapsto \Hom_{\ZZ}(M,\scrHom_E(X,N))
\]
transforms inductive limits in $X$ into projective limits of abelian groups, hence into projective limits of the underlying sets. It is therefore representable (IV, 1.4 and 1.2), and the representing object has an abelian sheaf structure. By definition, for every object $X$ of $E$ there is a canonical isomorphism
\begin{equation}
\scrHom_A(M,N)(X)=\Hom(X,\scrHom_A(M,N))\simeq
\Hom_A(M,\scrHom(X,N)).\tag{12.1.2}
\end{equation}
The isomorphism (12.1.1) follows from (12.1.2), the isomorphism $\scrHom(X,N)\simeq j_{X*}j_X^*N$ (10.8), and the adjunction formulas.

\subsection*{12.2}
The abelian sheaf $\scrHom_A(M,N)$ is called the \emph{sheaf of morphisms of $A$-modules from $M$ to $N$}. It is bifunctorial in $M$ and $N$. It follows from its definition and from 10.2 that it sends projective limits in the argument $N$ (respectively inductive limits in the argument $M$) to projective limits of abelian sheaves. In particular it is left exact in both arguments.

\begin{statement}{Proposition 12.3}
Let $(E,A)$ be a ringed topos, let $M,N$ be two right (respectively left) $A$-modules, and let $X$ be an object of $E$.
\begin{enumerate}[label=\alph*)]
\item There are canonical isomorphisms
\[
\scrHom_A(M,N)(X)\simeq \Hom_A(M,j_{X*}j_X^*N)
\simeq \Hom_{A|X}(j_X^*M,j_X^*N)
\simeq \Hom_A(j_{X!}j_X^*M,N).
\]
\item There is a canonical isomorphism
\[
\Phi_X:j_X^*\scrHom_A(M,N)\simeq \scrHom_{A|X}(j_X^*M,j_X^*N).
\]
Let moreover $P$ be a right (respectively left) $A|X$-module.
\item There is a canonical isomorphism
\[
j_{X*}\scrHom_{A|X}(j_X^*M,P)\simeq \scrHom_A(M,j_{X*}P).
\]
\item There is a canonical isomorphism
\[
\Xi_X:\scrHom_A(j_{X!}P,N)\simeq j_{X*}\scrHom_{A|X}(P,j_X^*N).
\]
\end{enumerate}
\end{statement}

The isomorphisms in a) follow from (12.1.2), from $\scrHom_E(X,N)\simeq j_{X*}j_X^*N$ (10.8), and from the adjunction formulas.

For b), for every object $Y$ of $E_{/X}$ one has a chain of isomorphisms
\begin{align*}
\Hom(Y,j_X^*\scrHom_A(M,N))
&\simeq \Hom(j_{X!}Y,\scrHom_A(M,N))\\
&\simeq \Hom_A(M,\scrHom_E(j_{X!}Y,N))\\
&\simeq \Hom_A(M,j_{X*}\scrHom_{E_{/X}}(Y,j_X^*N))\\
&\simeq \Hom_{A|X}(j_X^*M,\scrHom_{E_{/X}}(Y,j_X^*N))\\
&\simeq \Hom(Y,\scrHom_{A|X}(j_X^*M,j_X^*N)).
\end{align*}
For c), for every object $Y$ of $E$ one has
\begin{align*}
\Hom(Y,j_{X*}\scrHom_{A|X}(j_X^*M,P))
&\simeq \Hom(j_X^*Y,\scrHom_{A|X}(j_X^*M,P))\\
&\simeq \Hom_{A|X}(j_X^*M,\scrHom_{E_{/X}}(j_X^*Y,P))\\
&\simeq \Hom_A(M,j_{X*}\scrHom_{E_{/X}}(j_X^*Y,P))\\
&\simeq \Hom_A(M,\scrHom_E(Y,j_{X*}P))\\
&\simeq \Hom(Y,\scrHom_A(M,j_{X*}P)).
\end{align*}
For d), for every object $Y$ of $E$ one has
\begin{align*}
\Hom(Y,\scrHom_A(j_{X!}P,N))
&\simeq \Hom_A(j_{X!}P,\scrHom_E(Y,N))\\
&\simeq \Hom_{A|X}(P,j_X^*\scrHom_E(Y,N))\\
&\simeq \Hom_{A|X}(P,\scrHom_{E_{/X}}(j_X^*Y,j_X^*N))\\
&\simeq \Hom(j_X^*Y,\scrHom_{A|X}(P,j_X^*N))\\
&\simeq \Hom(Y,j_{X*}\scrHom_{A|X}(P,j_X^*N)).
\end{align*}

\begin{statement}{Corollary 12.4}
Let $C$ be a $\U$-site, let $A'$ be a $\U$-presheaf of rings on $C$, let $M$ be a $\U$-presheaf of left $A'$-modules, and let $N$ be a $\U$-sheaf of left $A'$-modules. Let $A$ denote the sheaf associated with $A'$, so that $N$ and $\underline aM$ are $A$-modules. Then the presheaf
\[
X\longmapsto \Hom_{A'|X}(j_X^*M,j_X^*N)
\]
is a $\U$-sheaf of abelian groups, canonically isomorphic to $\scrHom_A(\underline aM,N)$.
\end{statement}

Indeed, by III, 5.5 and III, 2.3, localization to $C_{/X}$ commutes with associated-sheaf functors (for $C$ and for $C_{/X}$). Hence one has isomorphisms, functorial in the variable object $X$ of $C$,
\begin{align*}
\Hom_{A'|X}(j_X^*M,j_X^*N)
&\simeq \Hom_{A|X}(\underline a\,j_X^*M,j_X^*N)\\
&\simeq \Hom_{A|X}(j_X^*\underline aM,j_X^*N)\\
&\simeq \scrHom_A(\underline aM,N)(X).
\end{align*}

\begin{statement}{Corollary 12.5}
Let $E$ be a topos, let $A,B,C$ be three sheaves of rings on $E$, let $M$ be an $A$-$B$-bimodule, and let $N$ be an $A$-$C$-bimodule. The abelian sheaf $\scrHom_A(M,N)$ has a canonical $B$-$C$-bimodule structure. In particular, when $A$ is a sheaf of commutative rings and when $M,N$ are $A$-modules, $\scrHom_A(M,N)$ has a canonical $A$-module structure. The canonical isomorphisms b), c), d) of 12.3 are isomorphisms of bimodules.
\end{statement}

We only give indications. For every object $X$ of $E$, one has
\[
\scrHom_A(M,N)(X)\simeq \Hom_{A|X}(j_X^*M,j_X^*N).
\]
Thus the abelian group $\scrHom_A(M,N)(X)$ has a canonical $B(X)$-$C(X)$-bimodule structure (11.1.3), which one checks is functorial in $X$. The other assertions are left to the reader.

\begin{statement}{Corollary 12.6}
Let $(E,A)$ be a ringed topos, let $X$ be an object of $E$, and let $N$ be an $A$-module, say on the left. There are canonical isomorphisms of $A$-modules
\[
\scrHom_E(X,N)\simeq j_{X*}j_X^*N\simeq \scrHom_A(A_X,N),
\]
where the $A$-module structure on $\scrHom_A(A_X,N)$ comes from the bimodule structure on $A_X$ (11.3.4).
\end{statement}

The first isomorphism follows from 10.7, and the second from the isomorphism of $A$-modules $N\simeq \scrHom_A(A,N)$ and from 12.3.

\begin{statement}{Proposition 12.7}
Let $(E,A)$ be a ringed topos, let $M$ be a right $A$-module and $N$ a left $A$-module. The functor which sends an abelian sheaf $P$ to the abelian group
\[
\Hom_A(M,\scrHom_{\underline{\ZZ}}(N,P))
\]
is representable by an abelian sheaf denoted $M\otimes_A N$ and called the tensor product of $M$ and $N$ over $A$.
\end{statement}

There is a canonical injection of $\Hom_A(M,\scrHom_{\underline{\ZZ}}(N,P))$ into
\[
\Hom(M,\scrHom(N,P))\simeq \Hom(M\times N,P).
\]
Returning to the definitions, one immediately sees that a morphism $f$ of sheaves of sets from $M\times N$ to $P$ comes from an element of $\Hom_A(M,\scrHom_{\ZZ}(N,P))$ if and only if, for every object $X$ of $E$, the map
\[
f(X):M(X)\times N(X)\longrightarrow P(X)
\]
is $A(X)$-bilinear. Call such morphisms $A$-bilinear morphisms. It remains to represent the functor of $A$-bilinear morphisms from $M\times N$ to $P$. One proceeds as in the ordinary case, \ie\ in the punctual topos. Consider the eight morphisms
\begin{align*}
&M\times M\times N\longrightarrow M\times N &&(1\le i\le 3),\\
&M\times N\times N\longrightarrow M\times N &&(4\le i\le 6),\\
&M\times A\times N\longrightarrow M\times N &&(7\le i\le 8),
\end{align*}
defined by
\begin{align*}
(1)&\quad (m_1,m_2,n)\longmapsto(m_1+m_2,n),&
(2)&\quad (m_1,m_2,n)\longmapsto(m_1,n),\\
(3)&\quad (m_1,m_2,n)\longmapsto(m_2,n),&
(4)&\quad (m,n_1,n_2)\longmapsto(m,n_1+n_2),\\
(5)&\quad (m,n_1,n_2)\longmapsto(m,n_1),&
(6)&\quad (m,n_1,n_2)\longmapsto(m,n_2),\\
(7)&\quad (m,a,n)\longmapsto(ma,n),&
(8)&\quad (m,a,n)\longmapsto(m,an).
\end{align*}
Here the formula $(i)$ describes the map $(i)(X)$ for variable objects $X$ of $E$. Let $L$ denote the ``free abelian sheaf generated by'' functor. It is clear that the largest quotient, in the category of abelian sheaves, of $L(M\times N)$ which equalizes $L(1)$ with $L(2)+L(3)$, $L(4)$ with $L(5)+L(6)$, and $L(7)$ with $L(8)$ represents the functor of $A$-bilinear morphisms from $M\times N$ to $P$.

\subsection*{12.8}
The functor
\[
P\longmapsto \Hom_A(N,\scrHom_{\underline{\ZZ}}(M,P))
\]
is also canonically isomorphic to the functor of $A$-bilinear maps from $M\times N$ to $P$. Hence $M\otimes_A N$ also represents the functor $\Hom_A(N,\scrHom_{\underline{\ZZ}}(M,P))$. Thus one has isomorphisms, functorial in all arguments,
\begin{align}
\Hom_{\underline{\ZZ}}(M\otimes_A N,P)&\simeq \Hom_A(M,\scrHom_{\underline{\ZZ}}(N,P)),\tag{12.8.1}\\
\Hom_{\underline{\ZZ}}(M\otimes_A N,P)&\simeq \Hom_A(N,\scrHom_{\underline{\ZZ}}(M,P)).\tag{12.8.2}
\end{align}

\subsection*{12.9}
It follows from 12.8, or from (12.8.1) and 12.2, that the functor $\otimes_A$ commutes with inductive limits in both arguments.

\begin{statement}{Proposition 12.10}
Let $C$ be a $\U$-site, let $A'$ be a presheaf of rings on $C$, let $M'$ (respectively $N'$) be a right (respectively left) $A'$-module, and let $A,M,N$ be the sheaves associated with $A',M',N'$, respectively. The sheaf associated with the presheaf
\[
X\longmapsto M'(X)\otimes_{A'(X)}N'(X),\qquad X\in\ob C,
\]
is canonically isomorphic to the sheaf $M\otimes_A N$.
\end{statement}

The presheaf $X\mapsto M'(X)\otimes_{A'(X)}N'(X)$ may be constructed from the presheaves $M'$, $N'$, and $A'$ by the operations indicated in the proof of 12.7. Since the associated-sheaf functor commutes with finite projective and inductive limits and with the free-abelian-object functor, formation of the tensor product commutes with the associated-sheaf functor.

\begin{statement}{Proposition 12.11}
Let $(E,A)$ be a ringed topos, let $M$ be a right $A$-module, let $N$ be a left $A$-module, and let $X$ be an object of $E$.
\begin{enumerate}[label=\alph*)]
\item There is a canonical isomorphism
\[
j_X^*(M\otimes_A N)\simeq (j_X^*M)\otimes_{A|X}(j_X^*N).
\]
Let moreover $P$ be a right $A|X$-module and $Q$ a left $A|X$-module.
\item There are canonical isomorphisms, the projection formulas,
\begin{align*}
j_{X!}(P\otimes_{A|X}j_X^*N)&\simeq (j_{X!}P)\otimes_A N,\\
j_{X!}(j_X^*M\otimes_{A|X}Q)&\simeq M\otimes_A(j_{X!}Q).
\end{align*}
\end{enumerate}
\end{statement}

For a), for every abelian sheaf $R$ on $E_{/X}$ there is a chain of isomorphisms, functorial in $R$,
\begin{align*}
\Hom_{\underline{\ZZ}}(j_X^*(M\otimes_A N),R)
&\simeq \Hom_{\underline{\ZZ}}(M\otimes_A N,j_{X*}R)\\
&\simeq \Hom_A(M,\scrHom_{\underline{\ZZ}}(N,j_{X*}R))\\
&\simeq \Hom_A(M,j_{X*}\scrHom_{\underline{\ZZ}}(j_X^*N,R))\\
&\simeq \Hom_{A|X}(j_X^*M,\scrHom_{\underline{\ZZ}}(j_X^*N,R))\\
&\simeq \Hom_{\underline{\ZZ}}(j_X^*M\otimes_{A|X}j_X^*N,R).
\end{align*}
For the first isomorphism in b), for every abelian sheaf $R$ there is a functorial chain
\begin{align*}
\Hom_{\underline{\ZZ}}(j_{X!}(P\otimes_{A|X}j_X^*N),R)
&\simeq \Hom_{\underline{\ZZ}}(P\otimes_{A|X}j_X^*N,j_X^*R)\\
&\simeq \Hom_{A|X}(P,\scrHom_{\underline{\ZZ}}(j_X^*N,j_X^*R))\\
&\simeq \Hom_{A|X}(P,j_X^*\scrHom_{\underline{\ZZ}}(N,R))\\
&\simeq \Hom_A(j_{X!}P,\scrHom_{\underline{\ZZ}}(N,R))\\
&\simeq \Hom_{\underline{\ZZ}}(j_{X!}P\otimes_A N,R).
\end{align*}
The second isomorphism is obtained in the same way using (12.8.2).

\begin{statement}{Corollary 12.12}
Let $E$ be a topos, let $A,B,C$ be three sheaves of rings on $E$, let $M$ be a $B$-$A$-bimodule, and let $N$ be an $A$-$C$-bimodule. The abelian sheaf $M\otimes_A N$ has a canonical $B$-$C$-bimodule structure. In particular, when $A$ is a sheaf of commutative rings and $M,N$ are $A$-modules, $M\otimes_A N$ has a canonical $A$-module structure. The isomorphisms of 12.9 are isomorphisms of bimodules.
\end{statement}

For every object $X$ of $E$, $j_X^*M$ is a right $A|X$-module on which the ring $B(X)$ acts on the left, and $j_X^*N$ is a left $A|X$-module on which $C(X)$ acts on the right. By functoriality, the abelian sheaf
\[
j_X^*M\otimes_{A|X}j_X^*N\simeq j_X^*(M\otimes_A N)
\]
has the structure of a $B(X)$-$C(X)$-object, functorially in $X$. Hence $M\otimes_A N$ has a $B$-$C$-bimodule structure. This $B(X)$-$C(X)$-structure on $j_X^*(M\otimes_A N)$ is reflected in the evident way on the functor of $A$-bilinear morphisms. One then checks that when $A$ is commutative and $M,N$ are $A$-modules, the right and left $A$-module structures so obtained are equal; hence $M\otimes_A N$ is an $A$-module. The last assertion is left to the reader.

\begin{statement}{Corollary 12.13}
Let $(E,A)$ be a ringed topos, let $M$ be a left (respectively right) $A$-module, and let $X$ be an object of $E$. With the notation $A_X$ of 11.3.3, there is a canonical isomorphism
\[
A_X\otimes_AM\simeq j_{X!}j_X^*M
\qquad
(\resp\ M\otimes_AA_X\simeq j_{X!}j_X^*M).
\]
\end{statement}

This follows from the projection formulas (12.11).

\begin{statement}{Proposition 12.14}
Let $E$ be a topos, let $A$ and $B$ be two sheaves of rings on $E$, let $M$ be a right $A$-module, let $N$ be an $A$-$B$-bimodule, and let $P$ be a right $B$-module. There is a canonical isomorphism
\[
\Hom_B(M\otimes_A N,P)\simeq \Hom_A(M,\scrHom_B(N,P)).
\]
\end{statement}

From (12.8.1) one obtains a canonical $A$-bilinear morphism
\[
\scrHom_{\ZZ}(N,P)\times N\longrightarrow P.
\]
Restricting to the subsheaf $\scrHom_B(N,P)\times N$, this gives an $A$-bilinear morphism
\[
\scrHom_B(N,P)\times N\longrightarrow P,
\]
which is immediately checked to be $B$-linear in the second factor. Thus one has a canonical morphism of $B$-modules
\[
\scrHom_B(N,P)\otimes_A N\longrightarrow P,
\]
and hence a canonical map, functorial in $M$,
\begin{equation}
\Hom_A(M,\scrHom_B(N,P))\longrightarrow \Hom_B(M\otimes_A N,P).\tag{12.14.1}
\end{equation}
We prove that this morphism of functors is an isomorphism. Since both sides send inductive limits in $M$ to projective limits, it is enough to prove that (12.14.1) is an isomorphism for $M$ running through a generating family of the category $E_A$. It is therefore enough to treat the case $M=A_X$, where $X$ is an object of $E$. The verification is then immediate.

\section{Morphisms of ringed topoi}

\begin{statement}{Definition 13.1}
Let $(E,A)$ and $(E',A')$ be two ringed topoi. A morphism of ringed topoi
\[
u:(E,A)\longrightarrow (E',A')
\]
is a pair $(m,\theta)$, where $m:E\to E'$ is a geometric morphism and $\theta:m^*A'\to A$ is a morphism of rings.
\end{statement}

\subsubsection*{13.1.1}
Since $m^*:E'\to E$ is left adjoint to $m_*:E\to E'$, giving a morphism of ringed topoi $(m,\theta):(E,A)\to(E',A')$ is equivalent to giving a geometric morphism $m:E\to E'$ and a morphism of sheaves of rings
\[
\Theta:A'\longrightarrow m_*A.
\]

\subsection*{13.2}
With a morphism of ringed topoi $u=(m,\theta):(E,A)\to(E',A')$ one associates two important functors between the corresponding categories of modules.

\subsubsection*{13.2.1. The direct-image functor for modules}
Let $M$ be a left (respectively right) $A$-module. The object $m_*M$ has a canonical $m_*A$-module structure; by restriction of scalars along the canonical morphism $\Theta:A'\to m_*A$, it becomes an $A'$-module, denoted $u_*M$ and called the \emph{direct image of $M$ by $u$}.

\subsubsection*{13.2.2. The inverse-image functor for modules}
Let $N$ be a left (respectively right) $A'$-module. The object $m^*N$ of $E$ has a canonical left (respectively right) $m^*A'$-module structure. The left $A$-module
\[
A\otimes_{m^*A'}m^*N
\]
(respectively the right $A$-module $m^*N\otimes_{m^*A'}A$), where $A$ is endowed with the $m^*A'$-module structure defined by $\theta:m^*A'\to A$, is denoted $u^*N$ and is called the \emph{inverse image of the module $N$ by the morphism of ringed topoi} $u$.

\subsubsection*{13.2.3}
For every $A$-module $M$, the underlying sheaf of sets and the underlying abelian sheaf of $u_*M$ are both $m_*M$. Thus the notation $u_*$ is most often also used for the direct image $m_*$ of sheaves of sets. On the other hand, for an $A'$-module $N$, the underlying sheaf of sets or abelian sheaf of $u^*N$ is in general neither equal nor isomorphic to $m^*N$ (except when $\theta$ is an isomorphism). One must therefore distinguish inverse image for modules from inverse image for sheaves of sets or abelian sheaves. The notation $u^{-1}:E'\to E$ is most often used for the functor $m^*$, the inverse image for sheaves of sets. The functor $u^{-1}$ is called the \emph{set-theoretic inverse image} by the morphism of ringed topoi $u$, as opposed to the inverse image $u^*$ ``in the sense of modules''.

\begin{statement}{Definition 13.3}
Let $(C,A)$ and $(C',A')$ be two $\U$-ringed sites. A morphism of ringed sites
\[
u:(C,A)\longrightarrow(C',A')
\]
is a pair $(m,\theta)$, where $m$ is a morphism from the site $C$ to the site $C'$ (IV, 4.9) and $\theta:m^*A'\to A$ is a morphism of rings.
\end{statement}

\subsubsection*{13.3.1}
A morphism of ringed topoi is a morphism of $\U$-ringed sites. A morphism of ringed sites gives rise to a morphism between the corresponding ringed topoi (11.1.1 and 4.9.1).

\begin{statement}{Proposition 13.4}
Let $u:(E,A)\to(E',A')$ be a morphism of ringed topoi.
\begin{enumerate}[label=\alph*)]
\item For a variable left (respectively right) $A$-module $M$, and for a variable left (respectively right) $A'$-module $N$, there are canonical bifunctorial isomorphisms, called adjunction isomorphisms,
\begin{align}
\Hom_A(u^*N,M)&\simeq \Hom_{A'}(N,u_*M),\tag{13.4.1}\\
u_*\scrHom_A(u^*N,M)&\simeq \scrHom_{A'}(N,u_*M).\tag{13.4.2}
\end{align}
\item For a variable object $X$ of $E'$, there is a canonical isomorphism, functorial in $X$,
\begin{equation}
u^*A'_X\longrightarrow A_{u^{-1}(X)},\tag{13.4.3}
\end{equation}
where $A'_X$ and $A_{u^{-1}(X)}$ denote the free modules generated by the indicated objects (11.3.3).
\item If $A'$ is commutative and if the canonical morphism $u^{-1}A'\to A$ is central (respectively if the canonical morphism $u^{-1}A'\to A$ is an isomorphism), then, for a variable right $A'$-module $M$ and a variable left $A'$-module $N$, there is a canonical bifunctorial isomorphism
\begin{align}
u^*M\otimes_Au^*N&\simeq u^*(M\otimes_{A'}N),\tag{13.4.4}\\
u^*M\otimes_Au^*N&\simeq u^{-1}(M\otimes_{A'}N)\qquad(\resp).\tag{13.4.5}
\end{align}
\end{enumerate}
\end{statement}

First there is a canonical isomorphism
\[
\Hom_A(u^*N,M)\simeq \Hom_{u^{-1}A'}(u^{-1}N,M)
\]
(12.14), and then an isomorphism
\[
\Hom_{u^{-1}A'}(u^{-1}N,M)\simeq \Hom_{A'}(N,u_*M)
\]
(III, 1.7); this proves (13.4.1). To exhibit (13.4.2), for every object $X$ of $E'$ one has functorial isomorphisms
\begin{align*}
\Hom_{E'}(X,u_*\scrHom_A(u^*N,M))
&\simeq \Hom_E(u^{-1}X,\scrHom_A(u^*N,M))\\
&\simeq \Hom_A(u^*N,\scrHom_E(u^{-1}X,M))\\
&\simeq \Hom_{A'}(N,u_*\scrHom_E(u^{-1}X,M))\\
&\simeq \Hom_{A'}(N,\scrHom_{E'}(X,u_*M))\\
&\simeq \Hom_{E'}(X,\scrHom_{A'}(N,u_*M)).
\end{align*}
Formula (13.4.5) is obtained from the corresponding adjunction formula and the definition of the tensor product (12.7). Formula (13.4.4) follows from (13.4.5), using commutativity of the tensor product when the base ring is commutative (12.8). Finally, to prove (13.4.3), consider the chain of isomorphisms
\begin{align*}
\Hom_A(u^*A'_X,M)&\simeq \Hom_{A'}(A'_X,u_*M) &&\text{by (13.4.1)},\\
&\simeq \Hom_{E'}(X,u_*M) &&\text{by (11.3.3)},\\
&\simeq \Hom_E(u^{-1}X,M) &&\text{by adjunction},\\
&\simeq \Hom_A(A_{u^{-1}X},M) &&\text{by (11.3.3).}
\end{align*}

\begin{statement}{Corollary 13.5}
Let $(E,A)$ be a ringed topos, let $x:P\to E$ be a point of $E$, and let $F\mapsto F_x$ be the associated fiber functor. The fiber functor at $x$ sends the tensor product of $A$-modules to the ordinary tensor product of $A_x$-modules.
\end{statement}

The tensor product in $P$ is the ordinary tensor product of modules, since $P$ is the punctual topos, namely the category of sets. The assertion follows from (13.4.4).

\begin{statement}{Corollary 13.6}
Let $u:(E,A)\to(E',A')$ be a morphism of ringed topoi. The direct-image functor $u_*$ for right or left modules commutes with projective limits and in particular is left exact. The inverse-image functor $u^*$ for right or left modules commutes with inductive limits and in particular is right exact.
\end{statement}

This follows from formula (13.4.1) (I, 2.11).

\subsection*{13.7}
Note that the inverse-image functor $u^*$ for modules is not, in general, exact, whereas the set-theoretic inverse-image functor $u^{-1}$ is exact. One can nevertheless assert the exactness of $u^*$ when the canonical morphism $u^{-1}A'\to A$ is flat on the right or on the left (V, 1.8). This is in particular the case when the canonical morphism $u^{-1}A'\to A$ is an isomorphism.

\subsection*{13.8}
Let $C$ be a $\U$-site, let $A'$ be a presheaf of rings on $C$, and let $A$ denote the associated sheaf. The category of presheaves of $A'$-modules which are sheaves is equivalent to the category of $A$-modules. One sometimes writes $\Hom_{A'}(M,N)$ for the group $\Hom_A(M,N)$ (11.1.2). Similarly, and abusively, one writes $\scrHom_{A'}(M,N)$ and $M\otimes_{A'}N$ for the sheaves $\scrHom_A(M,N)$ and $M\otimes_A N$. When $A'=k_E$ is the constant presheaf defined by an ordinary ring $k$, one also writes $\scrHom_k$ and $\otimes_k$ instead of $\scrHom_{k_E}$ and $\otimes_{k_E}$.

\begin{remarktext}{Exercise 13.9. Locally ringed topoi}
Let $(E,\Oo_E)$ be a commutatively ringed topos. For $X\in\ob E$ and $f\in\Oo_E(X)$, let $X_f$ be the largest subobject of $X$ on which $f$ is invertible.
\begin{enumerate}[label=\alph*)]
\item Show that for $X'\in\ob E_{/X}$, the section $f_{X'}$ is invertible if and only if the structural morphism $X'\to X$ factors through $X_f$, and that for $f,g\in\Oo_E(X)$ one has
\[
X_{fg}=X_f\cap X_g.
\]
\item Show that the following conditions (i)--(iii) are equivalent:
\begin{enumerate}[label=(\roman*)]
\item for $X\in\ob E$ and $f,g\in\Oo_E(X)$,
\[
X_{f+g}\subset \Sup(X_f,X_g);
\]
\item for $X\in\ob E$ and $f,g\in\Oo_E(X)$ such that $f+g$ is invertible, one has $X=\Sup(X_f,X_g)$;
\item for $X\in\ob E$ and $f\in\Oo_E(X)$, one has
\[
X=\Sup(X_f,X_{1-f}).
\]
\item Show that these conditions imply the following condition, and are equivalent to it if $E$ has enough points: for every point $p$ of $E$, the fiber ring $\Oo_{E,p}$ is a \emph{local ring}.
\end{enumerate}
When the equivalent conditions (i)--(iii) are satisfied, one says that $(E,\Oo_E)$ is a \emph{locally ringed topos}; similarly a ringed site is called a \emph{locally ringed site} if the corresponding ringed topos is locally ringed.
\item Let $(E,\Oo_E)$ and $(E',\Oo_{E'})$ be two locally ringed topoi, and let $f$ be a morphism of ringed topoi from the first to the second. Show that condition (i) below implies condition (ii), and is equivalent to it if $E$ has enough points:
\begin{enumerate}[label=(\roman*)]
\item for every $X'\in\ob E'$ and every $s'\in\Oo_{E'}(X')$, putting $X=f^{-1}(X')$ and $s=f^*(s')$, one has
\[
X'_{s'}=f^{-1}(X_s);
\]
\item for every point $p$ of $E$, putting $p'=f(p)$, the natural homomorphism on fibers
\[
\Oo_{E',p'}\longrightarrow \Oo_{E,p}
\]
is a \emph{local} homomorphism of local rings.
\end{enumerate}
When condition (i) is satisfied, one says that $f$ is a \emph{morphism of locally ringed topoi}, or, if confusion is possible, an \emph{admissible morphism of locally ringed topoi}. Denote by $\Hhomtoplocan(E,E')$ the full subcategory of the category $\Hhomtopan(E,E')$ of all morphisms of ringed topoi from $E$ to $E'$ defined by the admissible morphisms.
\item Suppose that $(E,\Oo_E)$ is the ringed topos associated with a ringed topological space $(X,\Oo_X)$ (IV, 2.1). Show that $(E,\Oo_E)$ is locally ringed if and only if $(X,\Oo_X)$ is locally ringed, \ie\ if and only if for every $x\in X$ the fiber $\Oo_{X,x}$ is a local ring. In criterion (iv) of b), it is enough to take $p$ in a conservative family of points of $E$. Let $f:(X,\Oo_X)\to(X',\Oo_{X'})$ be a morphism of ringed spaces, where $\Oo_X$ and $\Oo_{X'}$ are sheaves of local rings. Show that the corresponding morphism of ringed topoi
\[
\Top(X,\Oo_X)\longrightarrow \Top(X',\Oo_{X'})
\]
is admissible if and only if the same is true of $f$, \ie\ if and only if for every $x\in X$ the homomorphism $\Oo_{X',f(x)}\to\Oo_{X,x}$ is a \emph{local} homomorphism of local rings.
\item Let $(E,\Oo_E)$ and $(E',\Oo_{E'})$ be two locally ringed topoi, suppose that $E$ has enough points and that $(E',\Oo_{E'})$ is equivalent to the ringed topos defined by a \emph{scheme} $(X',\Oo_{X'})$. Prove that the category $\Hhomtoplocan(E,E')$ is equivalent to a \emph{discrete} category, \ie\ it is a rigid groupoid: every morphism is an isomorphism and every automorphism group is trivial.

Hint: reduce to the case where $(E,\Oo_E)$ is the punctual topos ringed by a field. Recall, by contrast, that $\Hhomtop(E,E')$ is not in general equivalent to a discrete category, even if $E$ and $E'$ are topoi defined by schemes, and even if $E$ is the punctual topos (cf. 4.2.3).
\item Let $E$ be a locally ringed topos, and let $P$ be the ringed topos defined by the ringed space $\Spec k$, where $k$ is a field. Show that admissible morphisms of locally ringed topoi $P\to E$ correspond to pairs $(p,u)$ consisting of a point $p$ of $E$ and an injection of $k(p)$ into $k$, where $k(p)$ is the residue field of the local ring $\Oo_{E,p}$. Generalize this to a statement exhibiting admissible morphisms from a punctual locally ringed topos into $E$ (generalizing EGA I, 2.4.4).

A \emph{geometric point of a locally ringed topos} $E$ means any admissible morphism into $E$ from the locally ringed topos defined by a ringed space of the form $\Spec(k)$, where $k$ is an \emph{algebraically closed} field. Define the \emph{category of geometric points of} $E$ (understood: corresponding to fields $k\in\U$), denoted $\Pptgeom(E)$, and a canonical functor
\[
\Pptgeom(E)\longrightarrow \Ppoint(E).
\]
Show that this functor is faithful if and only if $E$ has no point, \ie\ if and only if $\Ppoint(E)$ is the empty category.
\item Let $E$ be a $\U$-topos and let $\V$ be a universe such that $\U\in\V$. Define a 2-category $(\V\text{-}\U\text{-}\Top)_{/E}$ in terms of the object $E$ of the 2-category $(\V\text{-}\U\text{-}\Top)$, following the model of 5.14 a). When $E$ is ringed by a ring $A$, define similarly a 2-category $(\V\text{-}\U\text{-}\mathrm{Topan})_{/E}$; and when $E$ is locally ringed, using the notion of admissible morphism, define a 2-category $(\V\text{-}\U\text{-}\mathrm{Topanloc})_{/E}$, admitting $\Pptgeom(E)$ as a full subcategory.
\end{enumerate}
\end{remarktext}

\section{Modules on a topos defined by gluing}

\subsection*{14.1}
Let $(E,A)$ be a ringed topos, let $U$ be an open subtopos of $E$, let $F$ be its complementary closed subtopos, and let
\[
j:U\to E,\qquad i:F\to E
\]
be the canonical morphisms (9.3). The topos $U$ is ringed by $j^*A$ (11.2.1), and the morphism $i$, completed in the evident way, becomes a morphism of ringed topoi. In this section, the topos $F$ will be ringed by the sheaf $i^*A$, denoted $A|F$, and the morphism $j$, completed in the evident way, is then a morphism of ringed topoi.

\subsection*{14.2}
Thus one has two morphisms of ringed topoi
\[
j:(U,A|U)\longrightarrow(E,A),\qquad i:(F,A|F)\longrightarrow(E,A),
\]
and hence five functors between the corresponding module categories, written here for left modules:
\[
\xymatrix@C=2cm@R=1.3cm@M=.25cm{
{}_{A|U}U \ar@<14pt>[r]^{j_!}\ar@<-14pt>[r]_{j_*}&
{}_{A}E \ar[l]_{j^*}\ar@<7pt>[r]^{i^*}&
{}_{A|F}F\ar@<7pt>[l]_{i_*}
}
\tag{14.2.1}
\]
Each functor in (14.2.1) is left adjoint to the functor displayed below it.

\subsection*{14.3}
We know that the functor
\[
X\longmapsto (j^*X,i^*X;\\ i^*X\to i^*j_*j^*X)
\]
is an equivalence from $E$ to the gluing topos $(U,F,i^*j_*)$ (9.5.4). This equivalence induces an equivalence between the corresponding module categories. Hence the functor
\[
P\longmapsto (j^*P,i^*P;\\ i^*P\to i^*j_*j^*P)
\]
is an equivalence, denoted $\Phi$, from ${}_AE$ to the category $({}_{A|U}U,{}_{A|F}F,i^*j_*)$. The functors of (14.2.1), composed with $\Phi$ or $\Phi^{-1}$, are then the following functors:
\begin{align*}
\Phi\circ j_* &: M\longmapsto (M,i^*j_*M;\ i^*j_*M\xrightarrow{\id}i^*j_*M),\\
j^*\circ\Phi^{-1} &: (M,N;\ N\to i^*j_*M)\longmapsto M,\\
\Phi\circ j_! &: M\longmapsto (M,0;\ 0\to i^*j_*M),\\
\Phi\circ i_* &: N\longmapsto (0,N;\ N\to 0),\\
i^*\circ\Phi^{-1} &: (M,N;\ N\to i^*j_*M)\longmapsto N.
\end{align*}
As an exercise, the reader may spell out the various adjunction morphisms.

\subsection*{14.4}
Denote by
\begin{equation}
i^!:{}_AE\longrightarrow {}_{A|F}F\tag{14.4.1}
\end{equation}
the functor defined by the formula
\begin{equation}
i^!\circ\Phi^{-1}(M,N;N\xrightarrow{u}i^*j_*M)=\Ker(u).\tag{14.4.2}
\end{equation}

\begin{statement}{Proposition 14.5}
The functor $i^!$ is right adjoint to $i_*$. The adjunction morphism $i_*i^!\to\id$ is a monomorphism. The functors $i^!$, for variable rings, commute with restriction-of-scalars functors.
\end{statement}

It is clear that every morphism
\[
(0,N;0)\longrightarrow (M,N;N\xrightarrow{u}i^*j_*M)
\]
factors uniquely through $(0,\Ker(u);0)$, proving the adjunction property. The other assertions are trivial.

\begin{statement}{Proposition 14.6}
For every object $P$ of ${}_AE$ there are exact sequences, functorial in $P$,
\begin{gather}
0\longrightarrow j_!j^*P\longrightarrow P\longrightarrow i_*i^*P\longrightarrow 0,\tag{14.6.1}\\
0\longrightarrow i_*i^!P\longrightarrow P\longrightarrow j_*j^*P,\tag{14.6.2}
\end{gather}
where the non-trivial arrows are the adjunction morphisms.
\end{statement}

First observe that a sequence
\[
(M',N';u')\longrightarrow(M,N;u)\longrightarrow(M'',N'';u'')
\]
in $({}_{A|U}U,{}_{A|F}F;i^*j_*)$ is exact if and only if the corresponding sequences
\[
M'\to M\to M'',\qquad N'\to N\to N''
\]
are exact. The sequences (14.6.1) and (14.6.2) are transformed by $\Phi$ into sequences of the form
\begin{align*}
&0\longrightarrow(M,0;0)\longrightarrow(M,N;u)\longrightarrow(0,N;0)\longrightarrow0,\\
&0\longrightarrow(0,\Ker(u);0)\longrightarrow(M,N;u)\longrightarrow(M,i^*j_*M;\id).
\end{align*}
The exactness of these sequences is immediate.

\subsection*{14.7}
The exact sequence (14.6.2) gives a new interpretation of the functor $i_*i^!$. Let $X$ be an object of $E$. From (14.6.2) one obtains an exact sequence of abelian groups
\[
0\longrightarrow \Hom_E(X,i_*i^!P)\longrightarrow \Hom_E(X,P)\longrightarrow \Hom_E(X,j_*j^*P).
\]
Write again $U$ for the open of $E$, the final object of the open subtopos $U$. The adjunction properties of $j_*$, $j^*$, and $j_!^{\Set}$ imply that $\Hom_E(X,j_*j^*P)$ is canonically isomorphic to $\Hom_E(X\times U,P)$, and that the morphism
\[
\Hom_E(X,P)\longrightarrow \Hom_E(X,j_*j^*P)
\]
is the morphism defined by the monomorphism $X\times U\to X$ (IV, 5). In other words:

\begin{statement}{Proposition 14.8}
The sections of $i_*i^!P$ on $E_{/X}$ are the sections of $P$ on $E_{/X}$ whose support is contained in $F$ (notation of 9.3.5). Equivalently, $i_*i^!P$ is the largest subsheaf of $P$ with support contained in $F$.
\end{statement}

\subsection*{14.9}
By abuse of language, one sometimes says that $i_*i^!P$ is the submodule of $P$ defined by the sections of $P$ with support in $F$. By 8.5.3, this terminology merely extends to general topoi the terminology used for topoi of sheaves on topological spaces.

\begin{statement}{Proposition 14.10}
\begin{enumerate}[label=\arabic*)]
\item The functor $P\mapsto i_*i^*P$ is isomorphic to the functor
\[
P\mapsto i_*i^*A\otimes_A P.
\]
\item The functor $P\mapsto i_*i^!P$ is isomorphic to the functor
\[
P\mapsto \scrHom_A(i_*i^*A,P).
\]
\end{enumerate}
\end{statement}

The exact sequence (14.6.1), in the special case where $P$ is the $A$-module $A$, reads
\begin{equation}
0\longrightarrow A_U\longrightarrow A\longrightarrow i_*i^*A\longrightarrow0,\tag{14.10.1}
\end{equation}
where $A_U$ is the free $A$-module generated by the open $U$ corresponding to the open subtopos $U$ (9 and 11.3.3). For every $A$-module $P$, the $A$-modules $j_!j^*P$ and $j_*j^*P$ are canonically isomorphic to $A_U\otimes_AP$ and $\scrHom_A(A_U,P)$ respectively (12.3 and 12.6). Moreover, modulo these isomorphisms, the canonical morphisms $j_!j^*P\to P$ and $P\to j_*j^*P$ come from the monomorphism $A_U\to A$. From the exact sequence (14.10.1) one therefore obtains two exact sequences (12.2 and 12.11)
\begin{align*}
&j_!j^*P\longrightarrow P\longrightarrow i_*i^*A\otimes_AP\longrightarrow0,\\
&0\longrightarrow \scrHom_A(i_*i^*A,P)\longrightarrow P\longrightarrow j_*j^*P.
\end{align*}
Comparison with (14.6.1) and (14.6.2) gives the asserted isomorphisms.

\begin{thebibliography}{99}
\bibitem{1} M. Artin and B. Mazur, \emph{Homotopy of varieties in the etale topology}, in: Proceedings of a Conference on Local Fields, Driebergen 1966, Springer.
\bibitem{2} P. Gabriel and G. Zisman, \emph{Calculus of fractions and homotopy theory}, Ergebnisse der Mathematik, Bd. 35.
\bibitem{3} J. Giraud, \emph{Alg\`ebre homologique non commutative}, Grundlehren, Springer, 1971.
\bibitem{4} A. Grothendieck, \emph{Sur quelques points d'alg\`ebre homologique}, Tohoku Math. Journal, cited as [Toh].
\bibitem{5} A. Grothendieck, \emph{Fondements de la G\'eom\'etrie Alg\'ebrique}, Bourbaki expos\'es 1957/62, Secr\'etariat math\'ematique, 11 rue P. Curie, Paris.
\bibitem{6} A. Grothendieck, \emph{Crystals and the De Rham cohomology of schemes}, notes by I. Coates and O. Jussila, in: \emph{Dix expos\'es sur la cohomologie des sch\'emas}, North-Holland, 1969.
\bibitem{7} A. Grothendieck, \emph{Classes de Chern des repr\'esentations lin\'eaires des groupes discrets}, in: \emph{Dix expos\'es sur la cohomologie des sch\'emas}, North-Holland, 1969.
\bibitem{8} R. Godement, \emph{Th\'eorie des faisceaux}, Act. Scient. Ind. no. 1252, Hermann, Paris, 1958, cited as [TF].
\bibitem{9} M. Hakim, \emph{Topos annel\'es et sch\'emas relatifs}, multigraphed thesis, Orsay, 1967.
\bibitem{10} D. Mumford, \emph{Picard groups of moduli problems}, in \emph{Arithmetic Algebraic Geometry}, Harper's Series in Modern Mathematics.
\bibitem{11} Nguyen Dinh Ngoc, Th\`ese Sciences Math\'ematiques, Paris, 1963, no. 4995.
\bibitem{12} J. E. Roos, \emph{Distributivit\'e des $\varinjlim$ par rapport aux $\varprojlim$ des topos}, C. R. t.259 (1964), pp. 969--972, 1605--1608, 1801--1804.
\bibitem{13} P. Deligne and D. Mumford, \emph{The irreducibility of the space of curves of given genus}, Publ. Math. IHES no. 36 (1969).
\bibitem{14} B. Mitchell, \emph{Theory of categories}, Academic Press, 1965.
\bibitem{15} J. Lurie, \emph{Higher Topos Theory}, Annals of Mathematics Studies 170, Princeton University Press, 2009.
\bibitem{16} P. Johnstone, \emph{Topos theory}, London Mathematical Society Monographs 10, Academic Press, 1977.
\bibitem{17} F. Borceux and G. Van den Bossche, \emph{Structure des topologies d'un topos}, Cahiers de topologie et g\'eom\'etrie diff\'erentielle cat\'egoriques 25, no. 1 (1984), 37--39.
\end{thebibliography}

\end{document}
